%% ============================================================
%% sections/s355g.tex
%% United States Code, Title 21 - § 355g. Utilizing real world evidence
%% (FD&C Act § 505F). Source identifier: /us/usc/t21/s355g
%% Relevance to the Physical AI oncology trial context:
%% Utilizing real world evidence; the Cures Act RWE program relevant to evidence generation in Physical AI oncology trials.
%% Reproduced faithfully from the OLRC USLM XML (through Pub. L. 119-93).
%% No bracketed instructions or file names are inserted into the law.
%% ============================================================
\uscsection{§~355g.}{Utilizing real world evidence}
\uscprov{1}{\textbf{(a)}\textbf{ In general} The Secretary shall establish a program to evaluate the potential use of real world evidence-}
\uscprov{2}{\textbf{(1)} to help to support the approval of a new indication for a drug approved under section 355(c) of this title; and}
\uscprov{2}{\textbf{(2)} to help to support or satisfy postapproval study requirements.}
\uscprov{1}{\textbf{(b)}\textbf{ Real world evidence defined} In this section, the term ``real world evidence'' means data regarding the usage, or the potential benefits or risks, of a drug derived from sources other than traditional clinical trials.}
\uscprov{1}{\textbf{(c)}\textbf{ Program framework}}
\uscprov{2}{\textbf{(1)}\textbf{ In general} Not later than 2 years after December 13, 2016, the Secretary shall establish a draft framework for implementation of the program under this section.}
\uscprov{2}{\textbf{(2)}\textbf{ Contents of framework} The framework shall include information describing-}
\uscprov{3}{\textbf{(A)} the sources of real world evidence, including ongoing safety surveillance, observational studies, registries, claims, and patient-centered outcomes research activities;}
\uscprov{3}{\textbf{(B)} the gaps in data collection activities;}
\uscprov{3}{\textbf{(C)} the standards and methodologies for collection and analysis of real world evidence; and}
\uscprov{3}{\textbf{(D)} the priority areas, remaining challenges, and potential pilot opportunities that the program established under this section will address.}
\uscprov{2}{\textbf{(3)}\textbf{ Consultation}}
\uscprov{3}{\textbf{(A)}\textbf{ In general} In developing the program framework under this subsection, the Secretary shall consult with regulated industry, academia, medical professional organizations, representatives of patient advocacy organizations, consumer organizations, disease research foundations, and other interested parties.}
\uscprov{3}{\textbf{(B)}\textbf{ Process} The consultation under subparagraph (A) may be carried out through approaches such as-}
\uscprov{4}{\textbf{(i)} a public-private partnership with the entities described in such subparagraph in which the Secretary may participate;}
\uscprov{4}{\textbf{(ii)} a contract, grant, or other arrangement, as the Secretary determines appropriate, with such a partnership or an independent research organization; or}
\uscprov{4}{\textbf{(iii)} public workshops with the entities described in such subparagraph.}
\uscprov{1}{\textbf{(d)}\textbf{ Program implementation} The Secretary shall, not later than 3 years after December 13, 2016, and in accordance with the framework established under subsection (c), implement the program to evaluate the potential use of real world evidence.}
\uscprov{1}{\textbf{(e)}\textbf{ Guidance for industry} The Secretary shall-}
\uscprov{2}{\textbf{(1)} utilize the program established under subsection (a), its activities, and any subsequent pilots or written reports, to inform a guidance for industry on-}
\uscprov{3}{\textbf{(A)} the circumstances under which sponsors of drugs and the Secretary may rely on real world evidence for the purposes described in paragraphs (1) and (2) of subsection (a); and}
\uscprov{3}{\textbf{(B)} the appropriate standards and methodologies for collection and analysis of real world evidence submitted for such purposes;}
\uscprov{2}{\textbf{(2)} not later than 5 years after December 13, 2016, issue draft guidance for industry as described in paragraph (1); and}
\uscprov{2}{\textbf{(3)} not later than 18 months after the close of the public comment period for the draft guidance described in paragraph (2), issue revised draft guidance or final guidance.}
\uscprov{1}{\textbf{(f)}\textbf{ Rule of construction}}
\uscprov{2}{\textbf{(1)}\textbf{ In general} Subject to paragraph (2), nothing in this section prohibits the Secretary from using real world evidence for purposes not specified in this section, provided the Secretary determines that sufficient basis exists for any such nonspecified use.}
\uscprov{2}{\textbf{(2)}\textbf{ Standards of evidence and Secretary's authority} This section shall not be construed to alter-}
\uscprov{3}{\textbf{(A)} the standards of evidence under-}
\uscprov{4}{\textbf{(i)} subsection (c) or (d) of section 355 of this title, including the substantial evidence standard in such subsection (d); or}
\uscprov{4}{\textbf{(ii)} section 262(a) of title 42; or}
\uscprov{3}{\textbf{(B)} the Secretary's authority to require postapproval studies or clinical trials, or the standards of evidence under which studies or trials are evaluated.}
\uscsource{(June 25, 1938, ch. 675, §~505F, as added Pub. L. 114-255, div. A, title III, §~3022, Dec. 13, 2016, 130 Stat. 1096; amended Pub. L. 115-52, title IX, §~901(c), (d), Aug. 18, 2017, 131 Stat. 1076.)}
\usccrosshead{Editorial Notes}
\uscnotehead{Amendments}
\uscnotepar{2017-Subsec. (b). Pub. L. 115-52, §~901(c), substituted ``traditional'' for ``randomized''.}
\uscnotepar{Subsec. (d). Pub. L. 115-52, §~901(d), substituted ``3 years'' for ``2 years''.}

%% - DRAFTING INSTRUCTIONS (added for VVUQ-04 draft-bill; the reproduced statutory text above is unchanged) -
\begin{draftbox}{§ 355g (21 U.S.C. § 355g) - utilizing real world evidence}
\dstep{Role in the amendment: § 355g is the 21st Century Cures Act real-world-evidence program. It currently supports drug indications and postapproval study requirements. The conforming amendment extends the program to recognize verification, validation, and uncertainty quantification records and trial telemetry from Physical AI oncology investigations as a real-world data source, without disturbing the drug-approval core.}
\dstep{Amendatory action: Section 505F of the Federal Food, Drug, and Cosmetic Act (21 U.S.C. § 355g) is amended (1) in subsection (b), by inserting after "traditional clinical trials." the following: "Such term includes verification, validation, and uncertainty quantification records and trial telemetry generated under section 515D."; and (2) by adding at the end the following new subsection: "(g) Application to Physical AI oncology investigations. The Secretary shall, in carrying out the program under this section, evaluate the use of real world evidence, including records generated under section 515D, to support the evaluation of devices used in Physical AI oncology clinical investigations."}
\dstep{Substance and cross-references: keep the drug-approval purposes in subsections (a) through (f) intact; the new subsection (g) and the (b) clarification only add a recognized data source and a device-investigation use. Point only to section 515D for the verification records.}
\dstep{Source synthesis: process papers/VVUQ-04/instruct-bill/01-federal-statutory-framework.md (the Cures Act real-world-evidence program) and 09-vvuq-standards-clinical-trial-oncology.md (real-world data, decentralized-element trials, oncology evidence) together with papers/VVUQ-03/final-paper/sections/findings.tex (the evidentiary record) and papers/VVUQ-03/final-paper/sections/statutory\_text.tex. Resolve to references.bib keys pl-cures and usc-fdca.}
\dstep{Comparative print rendering: reproduce subsections (a) through (f) unchanged, show the inserted sentence in (b) wrapped in \cprinsert{...}, and show the new subsection (g) wrapped in \cprinsert{...} at the end.}
\dstep{Formatting: single hyphens only; § for every codified reference; even RaggedRight spacing; reword so no inserted clause ends in a one- or two-word line.}
\end{draftbox}
